% Options for packages loaded elsewhere
\PassOptionsToPackage{unicode}{hyperref}
\PassOptionsToPackage{hyphens}{url}
\PassOptionsToPackage{dvipsnames,svgnames,x11names}{xcolor}
%
\documentclass[
  letterpaper,
]{scrreprt}

\usepackage{amsmath,amssymb}
\usepackage{iftex}
\ifPDFTeX
  \usepackage[T1]{fontenc}
  \usepackage[utf8]{inputenc}
  \usepackage{textcomp} % provide euro and other symbols
\else % if luatex or xetex
  \usepackage{unicode-math}
  \defaultfontfeatures{Scale=MatchLowercase}
  \defaultfontfeatures[\rmfamily]{Ligatures=TeX,Scale=1}
\fi
\usepackage{lmodern}
\ifPDFTeX\else  
    % xetex/luatex font selection
  \setmonofont[]{JuliaMono}
\fi
% Use upquote if available, for straight quotes in verbatim environments
\IfFileExists{upquote.sty}{\usepackage{upquote}}{}
\IfFileExists{microtype.sty}{% use microtype if available
  \usepackage[]{microtype}
  \UseMicrotypeSet[protrusion]{basicmath} % disable protrusion for tt fonts
}{}
\makeatletter
\@ifundefined{KOMAClassName}{% if non-KOMA class
  \IfFileExists{parskip.sty}{%
    \usepackage{parskip}
  }{% else
    \setlength{\parindent}{0pt}
    \setlength{\parskip}{6pt plus 2pt minus 1pt}}
}{% if KOMA class
  \KOMAoptions{parskip=half}}
\makeatother
\usepackage{xcolor}
\setlength{\emergencystretch}{3em} % prevent overfull lines
\setcounter{secnumdepth}{0}
% Make \paragraph and \subparagraph free-standing
\ifx\paragraph\undefined\else
  \let\oldparagraph\paragraph
  \renewcommand{\paragraph}[1]{\oldparagraph{#1}\mbox{}}
\fi
\ifx\subparagraph\undefined\else
  \let\oldsubparagraph\subparagraph
  \renewcommand{\subparagraph}[1]{\oldsubparagraph{#1}\mbox{}}
\fi

\usepackage{color}
\usepackage{fancyvrb}
\newcommand{\VerbBar}{|}
\newcommand{\VERB}{\Verb[commandchars=\\\{\}]}
\DefineVerbatimEnvironment{Highlighting}{Verbatim}{commandchars=\\\{\}}
% Add ',fontsize=\small' for more characters per line
\usepackage{framed}
\definecolor{shadecolor}{RGB}{241,243,245}
\newenvironment{Shaded}{\begin{snugshade}}{\end{snugshade}}
\newcommand{\AlertTok}[1]{\textcolor[rgb]{0.68,0.00,0.00}{#1}}
\newcommand{\AnnotationTok}[1]{\textcolor[rgb]{0.37,0.37,0.37}{#1}}
\newcommand{\AttributeTok}[1]{\textcolor[rgb]{0.40,0.45,0.13}{#1}}
\newcommand{\BaseNTok}[1]{\textcolor[rgb]{0.68,0.00,0.00}{#1}}
\newcommand{\BuiltInTok}[1]{\textcolor[rgb]{0.00,0.23,0.31}{#1}}
\newcommand{\CharTok}[1]{\textcolor[rgb]{0.13,0.47,0.30}{#1}}
\newcommand{\CommentTok}[1]{\textcolor[rgb]{0.37,0.37,0.37}{#1}}
\newcommand{\CommentVarTok}[1]{\textcolor[rgb]{0.37,0.37,0.37}{\textit{#1}}}
\newcommand{\ConstantTok}[1]{\textcolor[rgb]{0.56,0.35,0.01}{#1}}
\newcommand{\ControlFlowTok}[1]{\textcolor[rgb]{0.00,0.23,0.31}{#1}}
\newcommand{\DataTypeTok}[1]{\textcolor[rgb]{0.68,0.00,0.00}{#1}}
\newcommand{\DecValTok}[1]{\textcolor[rgb]{0.68,0.00,0.00}{#1}}
\newcommand{\DocumentationTok}[1]{\textcolor[rgb]{0.37,0.37,0.37}{\textit{#1}}}
\newcommand{\ErrorTok}[1]{\textcolor[rgb]{0.68,0.00,0.00}{#1}}
\newcommand{\ExtensionTok}[1]{\textcolor[rgb]{0.00,0.23,0.31}{#1}}
\newcommand{\FloatTok}[1]{\textcolor[rgb]{0.68,0.00,0.00}{#1}}
\newcommand{\FunctionTok}[1]{\textcolor[rgb]{0.28,0.35,0.67}{#1}}
\newcommand{\ImportTok}[1]{\textcolor[rgb]{0.00,0.46,0.62}{#1}}
\newcommand{\InformationTok}[1]{\textcolor[rgb]{0.37,0.37,0.37}{#1}}
\newcommand{\KeywordTok}[1]{\textcolor[rgb]{0.00,0.23,0.31}{#1}}
\newcommand{\NormalTok}[1]{\textcolor[rgb]{0.00,0.23,0.31}{#1}}
\newcommand{\OperatorTok}[1]{\textcolor[rgb]{0.37,0.37,0.37}{#1}}
\newcommand{\OtherTok}[1]{\textcolor[rgb]{0.00,0.23,0.31}{#1}}
\newcommand{\PreprocessorTok}[1]{\textcolor[rgb]{0.68,0.00,0.00}{#1}}
\newcommand{\RegionMarkerTok}[1]{\textcolor[rgb]{0.00,0.23,0.31}{#1}}
\newcommand{\SpecialCharTok}[1]{\textcolor[rgb]{0.37,0.37,0.37}{#1}}
\newcommand{\SpecialStringTok}[1]{\textcolor[rgb]{0.13,0.47,0.30}{#1}}
\newcommand{\StringTok}[1]{\textcolor[rgb]{0.13,0.47,0.30}{#1}}
\newcommand{\VariableTok}[1]{\textcolor[rgb]{0.07,0.07,0.07}{#1}}
\newcommand{\VerbatimStringTok}[1]{\textcolor[rgb]{0.13,0.47,0.30}{#1}}
\newcommand{\WarningTok}[1]{\textcolor[rgb]{0.37,0.37,0.37}{\textit{#1}}}

\providecommand{\tightlist}{%
  \setlength{\itemsep}{0pt}\setlength{\parskip}{0pt}}\usepackage{longtable,booktabs,array}
\usepackage{calc} % for calculating minipage widths
% Correct order of tables after \paragraph or \subparagraph
\usepackage{etoolbox}
\makeatletter
\patchcmd\longtable{\par}{\if@noskipsec\mbox{}\fi\par}{}{}
\makeatother
% Allow footnotes in longtable head/foot
\IfFileExists{footnotehyper.sty}{\usepackage{footnotehyper}}{\usepackage{footnote}}
\makesavenoteenv{longtable}
\usepackage{graphicx}
\makeatletter
\def\maxwidth{\ifdim\Gin@nat@width>\linewidth\linewidth\else\Gin@nat@width\fi}
\def\maxheight{\ifdim\Gin@nat@height>\textheight\textheight\else\Gin@nat@height\fi}
\makeatother
% Scale images if necessary, so that they will not overflow the page
% margins by default, and it is still possible to overwrite the defaults
% using explicit options in \includegraphics[width, height, ...]{}
\setkeys{Gin}{width=\maxwidth,height=\maxheight,keepaspectratio}
% Set default figure placement to htbp
\makeatletter
\def\fps@figure{htbp}
\makeatother

\usepackage{mathrsfs}
\usepackage{twemojis}
\usepackage[normalem]{ulem}
\usepackage{scrlayer-scrpage}
\usepackage{textpos}
\automark{section}
\newcommand{\copyrightnotice}{\footnotesize\ Logan Johnson}
\newcommand{\longcopyrightnotice}{\footnotesize\ Logan Johnson 2023}
\setkomafont{pagefoot}{\upshape}
\ifoot*{\longcopyrightnotice}
\cfoot*{}
\ofoot*{\pagemark}

\makeatletter
\def\redsquiggly{\bgroup \markoverwith{\textcolor{red}{\lower3.5\p@\hbox{\sixly \char58}}}\ULon}
\def\brownsquiggly{\bgroup \markoverwith{\textcolor[HTML]{B8860B}{\lower3.5\p@\hbox{\sixly \char58}}}\ULon}
\def\bluesquiggly{\bgroup \markoverwith{\textcolor[HTML]{1E90FF}{\lower3.5\p@\hbox{\sixly \char58}}}\ULon}
\makeatother

\renewcommand{\NormalTok}[1]{\textcolor[HTML]{000000}{#1}}
\renewcommand{\KeywordTok}[1]{\textcolor[HTML]{0000FF}{#1}}
\renewcommand{\SpecialCharTok}[1]{}
\renewcommand{\ErrorTok}[1]{\redsquiggly{#1}}
\renewcommand{\WarningTok}[1]{\redsquiggly{\textcolor[HTML]{0000FF}{#1}}}
\renewcommand{\StringTok}[1]{\textcolor[HTML]{A52A2A}{#1}}
\renewcommand{\CommentTok}[1]{\textcolor[HTML]{008000}{#1}}
\renewcommand{\InformationTok}[1]{\textcolor[HTML]{D2691E}{\textbf{#1}}}
\renewcommand{\RegionMarkerTok}[1]{▼\:\textcolor[HTML]{008000}{\textbf{#1}}}
\renewcommand{\SpecialStringTok}[1]{\textcolor[HTML]{4682B4}{\textbf{#1}}}
\renewcommand{\ConstantTok}[1]{\textcolor[HTML]{DC143C}{#1}}
\renewcommand{\AnnotationTok}[1]{\brownsquiggly{#1}}
\renewcommand{\AlertTok}[1]{\brownsquiggly{\textcolor[HTML]{0000FF}{#1}}}
\renewcommand{\OtherTok}[1]{\bluesquiggly{#1}}
\renewcommand{\DocumentationTok}[1]{\bluesquiggly{\textcolor[HTML]{0000FF}{#1}}}

\newenvironment{ind}
	{\begin{list}{}{\setlength{\leftmargin}{1em}}\item\relax}
	{\end{list}}

%redefines Shaded so it can't break
\newcommand{\nobreakShaded}{\renewenvironment{Shaded}
	{\begin{tcolorbox}[frame hidden, enhanced, interior hidden, boxrule=0pt,
		borderline west={3pt}{0pt}{shadecolor}, sharp corners]}
	{\end{tcolorbox}}}

%Make end of environment ignore pars that come after it
\def\useignorespacesandallpars#1\ignorespaces\fi{%
#1\fi\ignorespacesandallpars}

\makeatletter
\def\ignorespacesandallpars{%
  \@ifnextchar\par
    {\expandafter\ignorespacesandallpars\@gobble}%
    {}%
}
\makeatother

\newenvironment{inpt}
	{\nobreakShaded\noindent\begin{minipage}[t]{0.63\textwidth}
		\uline{Lean File}}
	{\end{minipage}\hfill\useignorespacesandallpars}

\newenvironment{outpt}
	{\nobreakShaded\begin{minipage}[t]{0.32\textwidth}
		\uline{Tactic State in Infoview}}
	{\end{minipage}}

\newenvironment{bef}
	{\nobreakShaded\noindent\begin{minipage}[t]{0.475\textwidth}
		\uline{Tactic State Before Using Strategy}}
	{\end{minipage}\hfill\useignorespacesandallpars}

\newenvironment{aft}
	{\nobreakShaded\begin{minipage}[t]{0.475\textwidth}
		\uline{Tactic State After Using Strategy}}
	{\end{minipage}}

\newenvironment{numex}[1]
	{\begin{minipage}[t]{0.04\textwidth}\vspace{8pt}{#1}.
		\end{minipage}\nobreakShaded\begin{minipage}[t]{0.96\textwidth}\vspace{0pt}}
	{\end{minipage}}

\newenvironment{mdsk}
	{\medskip}
	{}

\newenvironment{absnobreak}
  {\par\nobreak\vfil\penalty0\vfilneg
   \vtop\bgroup}
  {\par\xdef\tpd{\the\prevdepth}\egroup
   \prevdepth=\tpd}

\newcommand{\excl}[1]{}
\newcommand{\incl}[1]{#1}

\newcommand{\setmin}{\mathbin{\backslash}}
\newcommand{\symmdiff}{\bigtriangleup}

\pagenumbering{roman}  %So front matter uses roman numerals.  Switch back to arabic at beginning of preface.
\publishers{\longcopyrightnotice}
\makeatletter
\makeatother
\makeatletter
\@ifpackageloaded{bookmark}{}{\usepackage{bookmark}}
\makeatother
\makeatletter
\@ifpackageloaded{caption}{}{\usepackage{caption}}
\AtBeginDocument{%
\ifdefined\contentsname
  \renewcommand*\contentsname{Table of contents}
\else
  \newcommand\contentsname{Table of contents}
\fi
\ifdefined\listfigurename
  \renewcommand*\listfigurename{List of Figures}
\else
  \newcommand\listfigurename{List of Figures}
\fi
\ifdefined\listtablename
  \renewcommand*\listtablename{List of Tables}
\else
  \newcommand\listtablename{List of Tables}
\fi
\ifdefined\figurename
  \renewcommand*\figurename{Figure}
\else
  \newcommand\figurename{Figure}
\fi
\ifdefined\tablename
  \renewcommand*\tablename{Table}
\else
  \newcommand\tablename{Table}
\fi
}
\@ifpackageloaded{float}{}{\usepackage{float}}
\floatstyle{ruled}
\@ifundefined{c@chapter}{\newfloat{codelisting}{h}{lop}}{\newfloat{codelisting}{h}{lop}[chapter]}
\floatname{codelisting}{Listing}
\newcommand*\listoflistings{\listof{codelisting}{List of Listings}}
\usepackage{amsthm}
\theoremstyle{remark}
\AtBeginDocument{\renewcommand*{\proofname}{Proof}}
\newtheorem*{remark}{Remark}
\newtheorem*{solution}{Solution}
\makeatother
\makeatletter
\@ifpackageloaded{caption}{}{\usepackage{caption}}
\@ifpackageloaded{subcaption}{}{\usepackage{subcaption}}
\makeatother
\makeatletter
\@ifpackageloaded{tcolorbox}{}{\usepackage[skins,breakable]{tcolorbox}}
\makeatother
\makeatletter
\@ifundefined{shadecolor}{\definecolor{shadecolor}{rgb}{.97, .97, .97}}
\makeatother
\makeatletter
\makeatother
\makeatletter
\makeatother
\ifLuaTeX
  \usepackage{selnolig}  % disable illegal ligatures
\fi
\IfFileExists{bookmark.sty}{\usepackage{bookmark}}{\usepackage{hyperref}}
\IfFileExists{xurl.sty}{\usepackage{xurl}}{} % add URL line breaks if available
\urlstyle{same} % disable monospaced font for URLs
\hypersetup{
  pdftitle={Untitled Lean Thesis},
  pdfauthor={Logan Johnson},
  colorlinks=true,
  linkcolor={blue},
  filecolor={Maroon},
  citecolor={Blue},
  urlcolor={Blue},
  pdfcreator={LaTeX via pandoc}}

\title{Untitled Lean Thesis}
\author{Logan Johnson}
\date{Invalid Date}

\begin{document}
\newcommand{\fromtop}[1]{%
  \dimexpr-1in-\topskip-\topmargin-\headheight-\headsep+#1\relax
}
\newcommand{\fromleft}[1]{%
  \dimexpr-1in-\oddsidemargin+#1\relax
}
\setlength{\textwidth}{6in}
\begin{titlepage}
\cleardoublepage
\thispagestyle{empty}
\begin{textblock*}{\paperwidth}(\fromleft{0cm},\fromtop{144pt})
\centering
UNTITLED LEAN THESIS
\end{textblock*}
\begin{textblock*}{\paperwidth}(\fromleft{0cm},\fromtop{300pt})
\centering
A Thesis Submitted to the Faculty of\\
\vspace{12pt}
Georgetown College\\
\vspace{12pt}
In Partial Fulfillment of the Requirements for the\\
\vspace{12pt}
Honors Program
\end{textblock*}
\begin{textblock*}{\paperwidth}(\fromleft{0cm},\fromtop{492pt})
\centering
By\\
\vspace{12pt}
Logan Johnson
\end{textblock*}
\begin{textblock*}{\paperwidth}(\fromleft{0cm},\fromtop{600pt})
\centering
Georgetown, Kentucky\\
\vspace{12pt}
May 2024
\end{textblock*}
\pagebreak
\thispagestyle{empty}
\begin{textblock*}{\paperwidth}(\fromleft{0cm},\fromtop{180pt})
\centering
Abstract\\
\vspace{12pt}
UNTITLED LEAN THESIS\\
\vspace{12pt}
Logan C. Johnson\\
\vspace{12pt}
idk dr burch?? maybe dr white???
\end{textblock*}
\begin{textblock*}{\textwidth}(\fromleft{1.5in},\fromtop{300pt})
\raggedright
Here is the text of your abstract. It goes on and on and on. It goes on like this for about
150 words, so it should all fit on this page. Note that the Abstract comes before the title page and
has no page number. The rest of this paragraph is a filler. It goes on like this for about 150
words, so it should all fit on this page. Note that the Abstract comes before the title page and has
no page number. It goes on like this for about 150 words, so it should all fit on this page. Note
that the Abstract comes before the title page and has no page number. It goes on like this for
about 150 words, so it should all fit on this page. Note that the Abstract comes before the title
page and has no page number. It goes on like this for about 150 words, so it should all fit on this
page. Note that the Abstract comes before the title page and has no page number. It goes on like
this for about 150 words, so it should all fit on this page. Note that the Abstract comes before the
title page and has no page number. It goes on like this for about 150 words, so it should all fit on
this page. Note that the Abstract comes before the title page and has no page number. If your
abstract is more than 250 words, consider shortening it.
\end{textblock*}
\pagebreak
\thispagestyle{empty}
\begin{textblock*}{\paperwidth}(\fromleft{2in},\fromtop{144pt})
APPROVED BY THE DIRECTOR OF HONORS THESES:
\end{textblock*}
\begin{textblock*}{\paperwidth}(\fromleft{3in},\fromtop{192pt})
\underline{\hspace{4.5in}}
\end{textblock*}
\begin{textblock*}{\paperwidth}(\fromleft{3in},\fromtop{216pt})
Dr. Homer White, Department of Mathematics
\end{textblock*}
\begin{textblock*}{\paperwidth}(\fromleft{2in},\fromtop{396pt})
APPROVED BY THE HONORS PROGRAM:
\end{textblock*}
\begin{textblock*}{\paperwidth}(\fromleft{2in},\fromtop{444pt})
\underline{\hspace{4.5in}}
\end{textblock*}
\begin{textblock*}{\paperwidth}(\fromleft{2in},\fromtop{468pt})
Dr. Barbara Burch, Director
\end{textblock*}
\begin{textblock*}{\paperwidth}(\fromleft{1.5in},\fromtop{612pt})
DATE:\underline{\hspace{2in}}
\end{textblock*}
\end{titlepage}%Can't be in preamble because Quarto loads amsthm too late.
\theoremstyle{plain}
\newtheorem*{thm}{Theorem}
\newcommand{\thmnm}{Theorem}
\newtheorem*{namedthm}{\thmnm}
\theoremstyle{definition}
\newtheorem*{dfn}{Definition}
\newcommand{\defnm}{Definition}
\newtheorem*{nameddfn}{\defnm}

\newenvironment{nthm}[1]
  {\renewcommand{\thmnm}{#1}\begin{namedthm}}
  {\end{namedthm}}

\newenvironment{ndfn}[1]
  {\renewcommand{\defnm}{#1}\begin{nameddfn}}
  {\end{nameddfn}}

\newenvironment{npf}[1]
  {\begin{proof}[#1]}
  {\end{proof}}

% Usage:
% ::: {.thm} will create an unnumbered thm environment with title Theorem.
% ::: {.nthm arguments="Name of Theorem"} will create an unnumbered theorem whose title is Name of Theorem.

\let\oldgreater\textgreater
\renewcommand{\textgreater}{\null\oldgreater}   % To prevent => changing to double arrow

\ifdefined\Shaded\renewenvironment{Shaded}{\begin{tcolorbox}[interior hidden, frame hidden, sharp corners, boxrule=0pt, enhanced, breakable, borderline west={3pt}{0pt}{shadecolor}]}{\end{tcolorbox}}\fi

\renewcommand*\contentsname{Table of contents}
{
\hypersetup{linkcolor=}
\setcounter{tocdepth}{1}
\tableofcontents
}
\bookmarksetup{startatroot}

\hypertarget{preface}{%
\chapter*{Preface}\label{preface}}
\addcontentsline{toc}{chapter}{Preface}

\markboth{Preface}{Preface}

\pagenumbering{arabic}
\markdouble{Preface}

I will not do my homework today.

\begin{Shaded}
\begin{Highlighting}[]
\FunctionTok{sum}\NormalTok{(}\DecValTok{4}\NormalTok{, }\DecValTok{7}\NormalTok{, }\DecValTok{3}\NormalTok{)}
\end{Highlighting}
\end{Shaded}

\begin{verbatim}
[1] 14
\end{verbatim}

Hello World

\hypertarget{making-chapters}{%
\section*{Making Chapters}\label{making-chapters}}
\addcontentsline{toc}{section}{Making Chapters}

\markright{Making Chapters}

I am using this section to figure out how to incorporate a table of
contents and different sections/chapters of the paper. This should prove
useful in the final thesis and allow readers to quickly jump to
important or interesting sections.

\hypertarget{incorporating-some-code}{%
\section*{Incorporating Some Code}\label{incorporating-some-code}}
\addcontentsline{toc}{section}{Incorporating Some Code}

\markright{Incorporating Some Code}

I will also be able to use some LaTeX equations within the document
which could halp to make the paper look quite nice.

If \(a < b\) and \(c \le d\), prove that \(a + c \le b + d\). It just so
happens that I was able to prove this using lean!

\begin{verbatim}
example (a b c d : ℝ) (h1: a < b) (h2 : c ≤ d) : a + c < b + d := by
  by_cases h3 : c = d
  rw [h3]
  apply add_lt_add_right h1
  push_neg at h3
  have h4 : c < d := by
    apply Ne.lt_of_le h3 h2
  apply add_lt_add h1 h4
  done
\end{verbatim}

Now to display the benefits of the lean infoview!

\hypertarget{showing-the-infoview-with-picture-sequences}{%
\section*{Showing the Infoview with Picture
Sequences}\label{showing-the-infoview-with-picture-sequences}}
\addcontentsline{toc}{section}{Showing the Infoview with Picture
Sequences}

\markright{Showing the Infoview with Picture Sequences}

Not going to do this now as I think the columns are far superior.

\hypertarget{showing-infoview-with-columns}{%
\section*{Showing Infoview with
Columns}\label{showing-infoview-with-columns}}
\addcontentsline{toc}{section}{Showing Infoview with Columns}

\markright{Showing Infoview with Columns}

As we can clearly see, this is the first step of the code and it can now
be explained with great ease. Now onto the next step!

\begin{inpt}

\begin{Shaded}
\begin{Highlighting}[]
\KeywordTok{theorem}\NormalTok{ Example\_3\_2\_4\_v2 (P Q R : }\KeywordTok{Prop}\NormalTok{)}
\NormalTok{    (h : P → (Q → R)) : ¬R → (P → ¬Q) := }\KeywordTok{by}
  \KeywordTok{assume}\NormalTok{ h2 : ¬R}
  \KeywordTok{assume}\NormalTok{ h3 : P}
  \SpecialCharTok{**}\WarningTok{done}\SpecialCharTok{::}
\end{Highlighting}
\end{Shaded}

\end{inpt}

\begin{outpt}

\begin{Shaded}
\begin{Highlighting}[]
\InformationTok{P Q R }\NormalTok{: Prop}
\InformationTok{h }\NormalTok{: P → Q → R}
\InformationTok{h2 }\NormalTok{: ¬R}
\InformationTok{h3 }\NormalTok{: P}
\NormalTok{⊢ ¬Q}
\end{Highlighting}
\end{Shaded}

\end{outpt}

Click on the \emph{Extensions} icon on the left side of the window,
which is circled in red in the image above. That will bring up a list of
available extensions:

\includegraphics{Images/FindExtension.png}

\hypertarget{acknowledgments}{%
\section*{Acknowledgments}\label{acknowledgments}}
\addcontentsline{toc}{section}{Acknowledgments}

\markright{Acknowledgments}

\bookmarksetup{startatroot}

\hypertarget{real-analysis}{%
\chapter{Real Analysis}\label{real-analysis}}

\begin{longtable}[]{@{}cc@{}}
\toprule\noalign{}
Symbol & Meaning \\
\midrule\noalign{}
\endhead
\bottomrule\noalign{}
\endlastfoot
\(\neg\) & not \\
\(\wedge\) & and \\
\(\vee\) & or \\
\(\to\) & if \ldots{} then \\
\(\leftrightarrow\) & iff (that is, if and only if) \\
\end{longtable}

\hypertarget{prop-laws}{}
\begin{longtable}[]{@{}
  >{\raggedright\arraybackslash}p{(\columnwidth - 6\tabcolsep) * \real{0.3226}}
  >{\centering\arraybackslash}p{(\columnwidth - 6\tabcolsep) * \real{0.2258}}
  >{\centering\arraybackslash}p{(\columnwidth - 6\tabcolsep) * \real{0.2258}}
  >{\centering\arraybackslash}p{(\columnwidth - 6\tabcolsep) * \real{0.2258}}@{}}
\toprule\noalign{}
\begin{minipage}[b]{\linewidth}\raggedright
Name
\end{minipage} & \begin{minipage}[b]{\linewidth}\centering
\end{minipage} & \begin{minipage}[b]{\linewidth}\centering
Equivalence
\end{minipage} & \begin{minipage}[b]{\linewidth}\centering
\end{minipage} \\
\midrule\noalign{}
\endhead
\bottomrule\noalign{}
\endlastfoot
De Morgan's Laws & \(\neg (P \wedge Q)\) & is equivalent to &
\(\neg P \vee \neg Q\) \\
& \(\neg (P \vee Q)\) & is equivalent to & \(\neg P \wedge \neg Q\) \\
Double Negation Law & \(\neg\neg P\) & is equivalent to & \(P\) \\
Conditional Laws & \(P \to Q\) & is equivalent to & \(\neg P \vee Q\) \\
& \(P \to Q\) & is equivalent to & \(\neg(P \wedge \neg Q)\) \\
Contrapositive Law & \(P \to Q\) & is equivalent to &
\(\neg Q \to \neg P\) \\
\end{longtable}

\begin{quote}
\(A \cap B = \{x \mid x \in A \wedge x \in B\} = {}\) the
\emph{intersection} of \(A\) and \(B\),

\(A \cup B = \{x \mid x \in A \vee x \in B\} = {}\) the \emph{union} of
\(A\) and \(B\),

\(A \setmin B = \{x \mid x \in A \wedge x \notin B\} = {}\) the
\emph{difference} of \(A\) and \(B\),

\(A \symmdiff B = (A \setmin B) \cup (B \setmin A) = {}\) the
\emph{symmetric difference} of \(A\) and \(B\).

\end{quote}

\bookmarksetup{startatroot}

\hypertarget{functional-programming}{%
\chapter{Functional Programming}\label{functional-programming}}

\begin{quote}
\(\forall x\,P(x)\) means ``for all \(x\), \(P(x)\),''

\end{quote}

\begin{longtable}[]{@{}ccc@{}}
\toprule\noalign{}
& Quantifier Negation Laws & \\
\midrule\noalign{}
\endhead
\bottomrule\noalign{}
\endlastfoot
\(\neg \exists x\,P(x)\) & is equivalent to &
\(\forall x\,\neg P(x)\) \\
\(\neg \forall x\,P(x)\) & is equivalent to &
\(\exists x\,\neg P(x)\) \\
\end{longtable}

\bookmarksetup{startatroot}

\hypertarget{lean-as-a-theorem-prover}{%
\chapter{Lean as a Theorem Prover}\label{lean-as-a-theorem-prover}}

\hypertarget{differences-from-paragraph-style-proofs}{%
\section{Differences From Paragraph Style
Proofs}\label{differences-from-paragraph-style-proofs}}

Despite the incredible power that lean could provide in the verification
of mathematical proofs, this does pose some difficulties, namely the
ease with which the aforementioned proofs can be written up. Typically,
proofs are simply written up in a paragraph style, where the steps being
taken and the theorems being applied are laid out in plain terms so that
it can be easily understood by fellow mathematicians. There are often
times when mathematicians will take things for granted or skip over
steps that they think the reader will either already know to be fact or
can easily reason out for themselves when writing out typical proofs.
This lax approach for conveying information simply does not work when
trying to communicate with technology, and a much more specific and
methodical approach must be adopted in order to take advantage of the
logical verification benefits. Thankfully, lean has a community working
to create libraries of previously proven theorems that can be applied to
speed up the writing and verification of future proofs. This thankfully
means that all proofs do not need to be taken all the way back to basic
axioms: Users can save time by avoiding proving adjacent theorems and
instead focus only on the immediately relevant steps of their proof.

For each of the following proofs, I will first provide a typical
``paragraph style'' version of the proof, so the differences between the
two can easily be compared. \#\# Inequality Addition

\begin{thm}
Whatever the problem is

\end{thm}

\begin{proof}

whatever the solution to that problem is. \excl{~□}\qedhere

\end{proof}

Step 1 Putting the initial theorem we want to show in lean code and
observing the final goal in the infoview.

\begin{inpt}

\begin{Shaded}
\begin{Highlighting}[]
\KeywordTok{example}\NormalTok{ (a b c d : ℝ) (h1: a \textless{} b) }
\NormalTok{    (h2 : c ≤ d) : a + c \textless{} b + d := }\KeywordTok{by}
  
  \SpecialCharTok{**}\WarningTok{done}\SpecialCharTok{::}
\end{Highlighting}
\end{Shaded}

\end{inpt}

\begin{outpt}

\begin{Shaded}
\begin{Highlighting}[]
\InformationTok{R}\NormalTok{: Type u\_1}
\InformationTok{inst✝}\NormalTok{: Ring R}
\InformationTok{abcd}\NormalTok{: ℝ}
\InformationTok{h1}\NormalTok{: a \textless{} b}
\InformationTok{h2}\NormalTok{: c ≤ d}
\NormalTok{⊢ a + c \textless{} b + d}
\end{Highlighting}
\end{Shaded}

\end{outpt}

Step 2

\begin{inpt}

\begin{Shaded}
\begin{Highlighting}[]
\KeywordTok{example}\NormalTok{ (a b c d : ℝ) (h1: a \textless{} b) }
\NormalTok{    (h2 : c ≤ d) : a + c \textless{} b + d := }\KeywordTok{by}
  \KeywordTok{by\_cases}\NormalTok{ h3 : c = d}

  \SpecialCharTok{**}\WarningTok{done}\SpecialCharTok{::}
\end{Highlighting}
\end{Shaded}

\end{inpt}

\begin{outpt}

\begin{Shaded}
\begin{Highlighting}[]
\InformationTok{R}\NormalTok{: Type u\_1}
\InformationTok{inst✝}\NormalTok{: Ring R}
\InformationTok{abcd}\NormalTok{: ℝ}
\InformationTok{h1}\NormalTok{: a \textless{} b}
\InformationTok{h2}\NormalTok{: c ≤ d}
\InformationTok{h3}\NormalTok{: c = d}
\NormalTok{⊢ a + c \textless{} b + d}
\end{Highlighting}
\end{Shaded}

\end{outpt}

Step 3

\begin{inpt}

\begin{Shaded}
\begin{Highlighting}[]
\KeywordTok{example}\NormalTok{ (a b c d : ℝ) (h1: a \textless{} b) }
\NormalTok{    (h2 : c ≤ d) : a + c \textless{} b + d := }\KeywordTok{by}
  \KeywordTok{by\_cases}\NormalTok{ h3 : c = d}
  \KeywordTok{rw}\NormalTok{ [h3]}

  \SpecialCharTok{**}\WarningTok{done}\SpecialCharTok{::}
\end{Highlighting}
\end{Shaded}

\end{inpt}

\begin{outpt}

\begin{Shaded}
\begin{Highlighting}[]
\InformationTok{R}\NormalTok{: Type u\_1}
\InformationTok{inst✝}\NormalTok{: Ring R}
\InformationTok{abcd}\NormalTok{: ℝ}
\InformationTok{h1}\NormalTok{: a \textless{} b}
\InformationTok{h2}\NormalTok{: c ≤ d}
\InformationTok{h3}\NormalTok{: c = d}
\NormalTok{⊢ a + d \textless{} b + d}
\end{Highlighting}
\end{Shaded}

\end{outpt}

Step 4

\begin{inpt}

\begin{Shaded}
\begin{Highlighting}[]
\KeywordTok{example}\NormalTok{ (a b c d : ℝ) (h1: a \textless{} b) }
\NormalTok{    (h2 : c ≤ d) : a + c \textless{} b + d := }\KeywordTok{by}
  \KeywordTok{by\_cases}\NormalTok{ h3 : c = d}
  \KeywordTok{rw}\NormalTok{ [h3]}
  \KeywordTok{apply}\NormalTok{ add\_lt\_add\_right h1}

  \SpecialCharTok{**}\WarningTok{done}\SpecialCharTok{::}
\end{Highlighting}
\end{Shaded}

\end{inpt}

\begin{outpt}

\begin{Shaded}
\begin{Highlighting}[]
\InformationTok{R}\NormalTok{: Type u\_1}
\InformationTok{inst✝}\NormalTok{: Ring R}
\InformationTok{abcd}\NormalTok{: ℝ}
\InformationTok{h1}\NormalTok{: a \textless{} b}
\InformationTok{h2}\NormalTok{: c ≤ d}
\InformationTok{h3}\NormalTok{: ¬c = d}
\NormalTok{⊢ a + c \textless{} b + d}
\end{Highlighting}
\end{Shaded}

\end{outpt}

Step 5

\begin{inpt}

\begin{Shaded}
\begin{Highlighting}[]
\KeywordTok{example}\NormalTok{ (a b c d : ℝ) (h1: a \textless{} b) }
\NormalTok{    (h2 : c ≤ d) : a + c \textless{} b + d := }\KeywordTok{by}
  \KeywordTok{by\_cases}\NormalTok{ h3 : c = d}
  \KeywordTok{rw}\NormalTok{ [h3]}
  \KeywordTok{apply}\NormalTok{ add\_lt\_add\_right h1}
  \KeywordTok{push\_neg} \KeywordTok{at}\NormalTok{ h3}

  \SpecialCharTok{**}\WarningTok{done}\SpecialCharTok{::}
\end{Highlighting}
\end{Shaded}

\end{inpt}

\begin{outpt}

\begin{Shaded}
\begin{Highlighting}[]
\InformationTok{R}\NormalTok{: Type u\_1}
\InformationTok{inst✝}\NormalTok{: Ring R}
\InformationTok{abcd}\NormalTok{: ℝ}
\InformationTok{h1}\NormalTok{: a \textless{} b}
\InformationTok{h2}\NormalTok{: c ≤ d}
\InformationTok{h3}\NormalTok{: c ≠ d}
\NormalTok{⊢ a + c \textless{} b + d}
\end{Highlighting}
\end{Shaded}

\end{outpt}

Step 6a We can see here that the lean infoview is now displaying my new
hypothesis as the current goal.

\begin{inpt}

\begin{Shaded}
\begin{Highlighting}[]
\KeywordTok{example}\NormalTok{ (a b c d : ℝ) (h1: a \textless{} b) }
\NormalTok{    (h2 : c ≤ d) : a + c \textless{} b + d := }\KeywordTok{by}
  \KeywordTok{by\_cases}\NormalTok{ h3 : c = d}
  \KeywordTok{rw}\NormalTok{ [h3]}
  \KeywordTok{apply}\NormalTok{ add\_lt\_add\_right h1}
  \KeywordTok{push\_neg} \KeywordTok{at}\NormalTok{ h3}
  \KeywordTok{have}\NormalTok{ h4 : c \textless{} d := }\KeywordTok{by}

  \SpecialCharTok{**}\WarningTok{done}\SpecialCharTok{::}
\end{Highlighting}
\end{Shaded}

\end{inpt}

\begin{outpt}

\begin{Shaded}
\begin{Highlighting}[]
\InformationTok{R}\NormalTok{: Type u\_1}
\InformationTok{inst✝}\NormalTok{: Ring R}
\InformationTok{abcd}\NormalTok{: ℝ}
\InformationTok{h1}\NormalTok{: a \textless{} b}
\InformationTok{h2}\NormalTok{: c ≤ d}
\InformationTok{h3}\NormalTok{: c ≠ d}
\NormalTok{⊢ c \textless{} d}
\end{Highlighting}
\end{Shaded}

\end{outpt}

Step 6b

\begin{inpt}

\begin{Shaded}
\begin{Highlighting}[]
\KeywordTok{example}\NormalTok{ (a b c d : ℝ) (h1: a \textless{} b) }
\NormalTok{    (h2 : c ≤ d) : a + c \textless{} b + d := }\KeywordTok{by}
  \KeywordTok{by\_cases}\NormalTok{ h3 : c = d}
  \KeywordTok{rw}\NormalTok{ [h3]}
  \KeywordTok{apply}\NormalTok{ add\_lt\_add\_right h1}
  \KeywordTok{push\_neg} \KeywordTok{at}\NormalTok{ h3}
  \KeywordTok{have}\NormalTok{ h4 : c \textless{} d := }\KeywordTok{by}
    \KeywordTok{apply}\NormalTok{ Ne.lt\_of\_le h3 h2}

  \SpecialCharTok{**}\WarningTok{done}\SpecialCharTok{::}
\end{Highlighting}
\end{Shaded}

\end{inpt}

\begin{outpt}

\begin{Shaded}
\begin{Highlighting}[]
\SpecialStringTok{No}\InformationTok{ }\SpecialStringTok{goals}
\end{Highlighting}
\end{Shaded}

\end{outpt}

Step 7

\begin{inpt}

\begin{Shaded}
\begin{Highlighting}[]
\KeywordTok{example}\NormalTok{ (a b c d : ℝ) (h1: a \textless{} b) }
\NormalTok{    (h2 : c ≤ d) : a + c \textless{} b + d := }\KeywordTok{by}
  \KeywordTok{by\_cases}\NormalTok{ h3 : c = d}
  \KeywordTok{rw}\NormalTok{ [h3]}
  \KeywordTok{apply}\NormalTok{ add\_lt\_add\_right h1}
  \KeywordTok{push\_neg} \KeywordTok{at}\NormalTok{ h3}
  \KeywordTok{have}\NormalTok{ h4 : c \textless{} d := }\KeywordTok{by}
    \KeywordTok{apply}\NormalTok{ Ne.lt\_of\_le h3 h2}
  \KeywordTok{apply}\NormalTok{ add\_lt\_add h1 h4}
  \KeywordTok{done}
\end{Highlighting}
\end{Shaded}

\end{inpt}

\begin{outpt}

\begin{Shaded}
\begin{Highlighting}[]
\SpecialStringTok{No}\InformationTok{ }\SpecialStringTok{goals}
\end{Highlighting}
\end{Shaded}

\end{outpt}

\hypertarget{a-is-less-than-or-equal-to-b}{%
\section{a is Less Than or Equal to
b}\label{a-is-less-than-or-equal-to-b}}

Setting up the problem

\begin{inpt}

\begin{Shaded}
\begin{Highlighting}[]
\KeywordTok{example}\NormalTok{ (a b : ℝ) (h1 : ∀ ε : ℝ, }
\NormalTok{    ε \textgreater{} 0 → a ≤ b + ε) : }
\NormalTok{    a ≤ b := }\KeywordTok{by}

  \SpecialCharTok{**}\WarningTok{done}\SpecialCharTok{::}
\end{Highlighting}
\end{Shaded}

\end{inpt}

\begin{outpt}

\begin{Shaded}
\begin{Highlighting}[]
\InformationTok{R}\NormalTok{: Type u\_1}
\InformationTok{inst✝}\NormalTok{: Ring R}
\InformationTok{ab}\NormalTok{: ℝ}
\InformationTok{h1}\NormalTok{: ∀ (ε : ℝ), ε \textgreater{} 0 → }
\SpecialCharTok{\textgreater{}\textgreater{}}\NormalTok{  a ≤ b + ε}
\NormalTok{⊢ a ≤ b}
\end{Highlighting}
\end{Shaded}

\end{outpt}

Step 1

\begin{inpt}

\begin{Shaded}
\begin{Highlighting}[]
\KeywordTok{example}\NormalTok{ (a b : ℝ) (h1 : ∀ ε : ℝ, }
\NormalTok{    ε \textgreater{} 0 → a ≤ b + ε) : }
\NormalTok{    a ≤ b := }\KeywordTok{by}
  \KeywordTok{by\_contra}\NormalTok{ h2}

  \SpecialCharTok{**}\WarningTok{done}\SpecialCharTok{::}
\end{Highlighting}
\end{Shaded}

\end{inpt}

\begin{outpt}

\begin{Shaded}
\begin{Highlighting}[]
\InformationTok{R}\NormalTok{: Type u\_1}
\InformationTok{inst✝}\NormalTok{: Ring R}
\InformationTok{ab}\NormalTok{: ℝ}
\InformationTok{h1}\NormalTok{: ∀ (ε : ℝ), ε \textgreater{} 0 → }
\SpecialCharTok{\textgreater{}\textgreater{}}\NormalTok{  a ≤ b + ε}
\InformationTok{h2}\NormalTok{: ¬a ≤ b}
\NormalTok{⊢ False}
\end{Highlighting}
\end{Shaded}

\end{outpt}

Step 2

\begin{inpt}

\begin{Shaded}
\begin{Highlighting}[]
\KeywordTok{example}\NormalTok{ (a b : ℝ) (h1 : ∀ ε : ℝ, }
\NormalTok{    ε \textgreater{} 0 → a ≤ b + ε) : }
\NormalTok{    a ≤ b := }\KeywordTok{by}
  \KeywordTok{by\_contra}\NormalTok{ h2}
  \KeywordTok{push\_neg} \KeywordTok{at}\NormalTok{ h2}
  \KeywordTok{let}\NormalTok{ ε := (a {-} b) / 2}

  \SpecialCharTok{**}\WarningTok{done}\SpecialCharTok{::}
\end{Highlighting}
\end{Shaded}

\end{inpt}

\begin{outpt}

\begin{Shaded}
\begin{Highlighting}[]
\InformationTok{R}\NormalTok{: Type u\_1}
\InformationTok{inst✝}\NormalTok{: Ring R}
\InformationTok{ab}\NormalTok{: ℝ}
\InformationTok{h1}\NormalTok{: ∀ (ε : ℝ), ε \textgreater{} 0 → }
\SpecialCharTok{\textgreater{}\textgreater{}}\NormalTok{  a ≤ b + ε}
\InformationTok{h2}\NormalTok{: b \textless{} a}
\InformationTok{ε}\NormalTok{: ℝ := (a {-} b) / 2}
\NormalTok{⊢ False}
\end{Highlighting}
\end{Shaded}

\end{outpt}

Step 3

\begin{inpt}

\begin{Shaded}
\begin{Highlighting}[]
\KeywordTok{example}\NormalTok{ (a b : ℝ) (h1 : ∀ ε : ℝ, }
\NormalTok{    ε \textgreater{} 0 → a ≤ b + ε) : }
\NormalTok{    a ≤ b := }\KeywordTok{by}
  \KeywordTok{by\_contra}\NormalTok{ h2}
  \KeywordTok{push\_neg} \KeywordTok{at}\NormalTok{ h2}
  \KeywordTok{let}\NormalTok{ ε := (a {-} b) / 2}
  \KeywordTok{have}\NormalTok{ h3 : ε \textgreater{} 0 := }\KeywordTok{by}

    \SpecialCharTok{**}\WarningTok{done}\SpecialCharTok{::}

  \SpecialCharTok{**}\WarningTok{done}\SpecialCharTok{::}
\end{Highlighting}
\end{Shaded}

\end{inpt}

\begin{outpt}

\begin{Shaded}
\begin{Highlighting}[]
\InformationTok{R}\NormalTok{: Type u\_1}
\InformationTok{inst✝}\NormalTok{: Ring R}
\InformationTok{ab}\NormalTok{: ℝ}
\InformationTok{h1}\NormalTok{: ∀ (ε : ℝ), ε \textgreater{} 0 → }
\SpecialCharTok{\textgreater{}\textgreater{}}\NormalTok{  a ≤ b + ε}
\InformationTok{h2}\NormalTok{: b \textless{} a}
\InformationTok{ε}\NormalTok{: ℝ := (a {-} b) / 2}
\NormalTok{⊢ ε \textgreater{} 0}
\end{Highlighting}
\end{Shaded}

\end{outpt}

Step 4

\begin{inpt}

\begin{Shaded}
\begin{Highlighting}[]
\KeywordTok{example}\NormalTok{ (a b : ℝ) (h1 : ∀ ε : ℝ, }
\NormalTok{    ε \textgreater{} 0 → a ≤ b + ε) : }
\NormalTok{    a ≤ b := }\KeywordTok{by}
  \KeywordTok{by\_contra}\NormalTok{ h2}
  \KeywordTok{push\_neg} \KeywordTok{at}\NormalTok{ h2}
  \KeywordTok{let}\NormalTok{ ε := (a {-} b) / 2}
  \KeywordTok{have}\NormalTok{ h3 : ε \textgreater{} 0 := }\KeywordTok{by}
\NormalTok{    refine half\_pos ?h}

      \SpecialCharTok{**}\WarningTok{done}\SpecialCharTok{::}

  \SpecialCharTok{**}\WarningTok{done}\SpecialCharTok{::}
\end{Highlighting}
\end{Shaded}

\end{inpt}

\begin{outpt}

\begin{Shaded}
\begin{Highlighting}[]
\InformationTok{R}\NormalTok{: Type u\_1}
\InformationTok{inst✝}\NormalTok{: Ring R}
\InformationTok{ab}\NormalTok{: ℝ}
\InformationTok{h1}\NormalTok{: ∀ (ε : ℝ), ε \textgreater{} 0 → }
\SpecialCharTok{\textgreater{}\textgreater{}}\NormalTok{  a ≤ b + ε}
\InformationTok{h2}\NormalTok{: b \textless{} a}
\InformationTok{ε}\NormalTok{: ℝ := (a {-} b) / 2}
\NormalTok{⊢ 0 \textless{} a {-} b}
\end{Highlighting}
\end{Shaded}

\end{outpt}

Step 5

\begin{inpt}

\begin{Shaded}
\begin{Highlighting}[]
\KeywordTok{example}\NormalTok{ (a b : ℝ) (h1 : ∀ ε : ℝ, }
\NormalTok{    ε \textgreater{} 0 → a ≤ b + ε) : }
\NormalTok{    a ≤ b := }\KeywordTok{by}
  \KeywordTok{by\_contra}\NormalTok{ h2}
  \KeywordTok{push\_neg} \KeywordTok{at}\NormalTok{ h2}
  \KeywordTok{let}\NormalTok{ ε := (a {-} b) / 2}
  \KeywordTok{have}\NormalTok{ h3 : ε \textgreater{} 0 := }\KeywordTok{by}
\NormalTok{    refine half\_pos ?h}
    \KeywordTok{exact}\NormalTok{ Iff.mpr sub\_pos h2}
    \KeywordTok{done}

  \SpecialCharTok{**}\WarningTok{done}\SpecialCharTok{::}
\end{Highlighting}
\end{Shaded}

\end{inpt}

\begin{outpt}

\begin{Shaded}
\begin{Highlighting}[]
\InformationTok{R}\NormalTok{: Type u\_1}
\InformationTok{inst✝}\NormalTok{: Ring R}
\InformationTok{ab}\NormalTok{: ℝ}
\InformationTok{h1}\NormalTok{: ∀ (ε : ℝ), ε \textgreater{} 0 → }
\SpecialCharTok{\textgreater{}\textgreater{}}\NormalTok{  a ≤ b + ε}
\InformationTok{h2}\NormalTok{: b \textless{} a}
\InformationTok{ε}\NormalTok{: ℝ := (a {-} b) / 2}
\InformationTok{h3}\NormalTok{: ε \textgreater{} 0}
\NormalTok{⊢ False}
\end{Highlighting}
\end{Shaded}

\end{outpt}

Step 6

\begin{inpt}

\begin{Shaded}
\begin{Highlighting}[]
\KeywordTok{example}\NormalTok{ (a b : ℝ) (h1 : ∀ ε : ℝ, }
\NormalTok{    ε \textgreater{} 0 → a ≤ b + ε) : }
\NormalTok{    a ≤ b := }\KeywordTok{by}
  \KeywordTok{by\_contra}\NormalTok{ h2}
  \KeywordTok{push\_neg} \KeywordTok{at}\NormalTok{ h2}
  \KeywordTok{let}\NormalTok{ ε := (a {-} b) / 2}
  \KeywordTok{have}\NormalTok{ h3 : ε \textgreater{} 0 := }\KeywordTok{by}
\NormalTok{    refine half\_pos ?h}
    \KeywordTok{exact}\NormalTok{ Iff.mpr sub\_pos h2}
    \KeywordTok{done}
  \KeywordTok{have}\NormalTok{ h4 : a ≤ b + ε := }\KeywordTok{by}

    \SpecialCharTok{**}\WarningTok{done}\SpecialCharTok{::}

  \SpecialCharTok{**}\WarningTok{done}\SpecialCharTok{::}
\end{Highlighting}
\end{Shaded}

\end{inpt}

\begin{outpt}

\begin{Shaded}
\begin{Highlighting}[]
\InformationTok{R}\NormalTok{: Type u\_1}
\InformationTok{inst✝}\NormalTok{: Ring R}
\InformationTok{ab}\NormalTok{: ℝ}
\InformationTok{h1}\NormalTok{: ∀ (ε : ℝ), ε \textgreater{} 0 → }
\SpecialCharTok{\textgreater{}\textgreater{}}\NormalTok{  a ≤ b + ε}
\InformationTok{h2}\NormalTok{: b \textless{} a}
\InformationTok{ε}\NormalTok{: ℝ := (a {-} b) / 2}
\InformationTok{h3}\NormalTok{: ε \textgreater{} 0}
\NormalTok{⊢ a ≤ b + ε}
\end{Highlighting}
\end{Shaded}

\end{outpt}

Step 7

\begin{inpt}

\begin{Shaded}
\begin{Highlighting}[]
\KeywordTok{example}\NormalTok{ (a b : ℝ) (h1 : ∀ ε : ℝ, }
\NormalTok{    ε \textgreater{} 0 → a ≤ b + ε) : }
\NormalTok{    a ≤ b := }\KeywordTok{by}
  \KeywordTok{by\_contra}\NormalTok{ h2}
  \KeywordTok{push\_neg} \KeywordTok{at}\NormalTok{ h2}
  \KeywordTok{let}\NormalTok{ ε := (a {-} b) / 2}
  \KeywordTok{have}\NormalTok{ h3 : ε \textgreater{} 0 := }\KeywordTok{by}
\NormalTok{    refine half\_pos ?h}
    \KeywordTok{exact}\NormalTok{ Iff.mpr sub\_pos h2}
    \KeywordTok{done}
  \KeywordTok{have}\NormalTok{ h4 : a ≤ b + ε := }\KeywordTok{by}
    \KeywordTok{apply}\NormalTok{ h1}

    \SpecialCharTok{**}\WarningTok{done}\SpecialCharTok{::}

  \SpecialCharTok{**}\WarningTok{done}\SpecialCharTok{::}
\end{Highlighting}
\end{Shaded}

\end{inpt}

\begin{outpt}

\begin{Shaded}
\begin{Highlighting}[]
\InformationTok{R}\NormalTok{: Type u\_1}
\InformationTok{inst✝}\NormalTok{: Ring R}
\InformationTok{ab}\NormalTok{: ℝ}
\InformationTok{h1}\NormalTok{: ∀ (ε : ℝ), ε \textgreater{} 0 → }
\SpecialCharTok{\textgreater{}\textgreater{}}\NormalTok{  a ≤ b + ε}
\InformationTok{h2}\NormalTok{: b \textless{} a}
\InformationTok{ε}\NormalTok{: ℝ := (a {-} b) / 2}
\InformationTok{h3}\NormalTok{: ε \textgreater{} 0}
\NormalTok{⊢ ε \textgreater{} 0}
\end{Highlighting}
\end{Shaded}

\end{outpt}

Step 8

\begin{inpt}

\begin{Shaded}
\begin{Highlighting}[]
\KeywordTok{example}\NormalTok{ (a b : ℝ) (h1 : ∀ ε : ℝ, }
\NormalTok{    ε \textgreater{} 0 → a ≤ b + ε) : }
\NormalTok{    a ≤ b := }\KeywordTok{by}
  \KeywordTok{by\_contra}\NormalTok{ h2}
  \KeywordTok{push\_neg} \KeywordTok{at}\NormalTok{ h2}
  \KeywordTok{let}\NormalTok{ ε := (a {-} b) / 2}
  \KeywordTok{have}\NormalTok{ h3 : ε \textgreater{} 0 := }\KeywordTok{by}
\NormalTok{    refine half\_pos ?h}
    \KeywordTok{exact}\NormalTok{ Iff.mpr sub\_pos h2}
    \KeywordTok{done}
  \KeywordTok{have}\NormalTok{ h4 : a ≤ b + ε := }\KeywordTok{by}
    \KeywordTok{apply}\NormalTok{ h1}
    \KeywordTok{apply}\NormalTok{ h3}
    \KeywordTok{done}

  \SpecialCharTok{**}\WarningTok{done}\SpecialCharTok{::}
\end{Highlighting}
\end{Shaded}

\end{inpt}

\begin{outpt}

\begin{Shaded}
\begin{Highlighting}[]
\InformationTok{R}\NormalTok{: Type u\_1}
\InformationTok{inst✝}\NormalTok{: Ring R}
\InformationTok{ab}\NormalTok{: ℝ}
\InformationTok{h1}\NormalTok{: ∀ (ε : ℝ), ε \textgreater{} 0 → }
\SpecialCharTok{\textgreater{}\textgreater{}}\NormalTok{  a ≤ b + ε}
\InformationTok{h2}\NormalTok{: b \textless{} a}
\InformationTok{ε}\NormalTok{: ℝ := (a {-} b) / 2}
\InformationTok{h3}\NormalTok{: ε \textgreater{} 0}
\InformationTok{h4}\NormalTok{: a ≤ b + ε}
\NormalTok{⊢ False}
\end{Highlighting}
\end{Shaded}

\end{outpt}

Step 9

\begin{inpt}

\begin{Shaded}
\begin{Highlighting}[]
\KeywordTok{example}\NormalTok{ (a b : ℝ) (h1 : ∀ ε : ℝ, }
\NormalTok{    ε \textgreater{} 0 → a ≤ b + ε) : }
\NormalTok{    a ≤ b := }\KeywordTok{by}
  \KeywordTok{by\_contra}\NormalTok{ h2}
  \KeywordTok{push\_neg} \KeywordTok{at}\NormalTok{ h2}
  \KeywordTok{let}\NormalTok{ ε := (a {-} b) / 2}
  \KeywordTok{have}\NormalTok{ h3 : ε \textgreater{} 0 := }\KeywordTok{by}
\NormalTok{    refine half\_pos ?h}
    \KeywordTok{exact}\NormalTok{ Iff.mpr sub\_pos h2}
    \KeywordTok{done}
  \KeywordTok{have}\NormalTok{ h4 : a ≤ b + ε := }\KeywordTok{by}
    \KeywordTok{apply}\NormalTok{ h1}
    \KeywordTok{apply}\NormalTok{ h3}
    \KeywordTok{done}
\NormalTok{  dsimp }\KeywordTok{at}\NormalTok{ h4}

  \SpecialCharTok{**}\WarningTok{done}\SpecialCharTok{::}
\end{Highlighting}
\end{Shaded}

\end{inpt}

\begin{outpt}

\begin{Shaded}
\begin{Highlighting}[]
\InformationTok{R}\NormalTok{: Type u\_1}
\InformationTok{inst✝}\NormalTok{: Ring R}
\InformationTok{ab}\NormalTok{: ℝ}
\InformationTok{h1}\NormalTok{: ∀ (ε : ℝ), ε \textgreater{} 0 → }
\SpecialCharTok{\textgreater{}\textgreater{}}\NormalTok{  a ≤ b + ε}
\InformationTok{h2}\NormalTok{: b \textless{} a}
\InformationTok{ε}\NormalTok{: ℝ := (a {-} b) / 2}
\InformationTok{h3}\NormalTok{: ε \textgreater{} 0}
\InformationTok{h4}\NormalTok{: a ≤ b + (a {-} b) / 2}
\NormalTok{⊢ False}
\end{Highlighting}
\end{Shaded}

\end{outpt}

Step 10

\begin{inpt}

\begin{Shaded}
\begin{Highlighting}[]
\KeywordTok{example}\NormalTok{ (a b : ℝ) (h1 : ∀ ε : ℝ, }
\NormalTok{    ε \textgreater{} 0 → a ≤ b + ε) : }
\NormalTok{    a ≤ b := }\KeywordTok{by}
  \KeywordTok{by\_contra}\NormalTok{ h2}
  \KeywordTok{push\_neg} \KeywordTok{at}\NormalTok{ h2}
  \KeywordTok{let}\NormalTok{ ε := (a {-} b) / 2}
  \KeywordTok{have}\NormalTok{ h3 : ε \textgreater{} 0 := }\KeywordTok{by}
\NormalTok{    refine half\_pos ?h}
    \KeywordTok{exact}\NormalTok{ Iff.mpr sub\_pos h2}
    \KeywordTok{done}
  \KeywordTok{have}\NormalTok{ h4 : a ≤ b + ε := }\KeywordTok{by}
    \KeywordTok{apply}\NormalTok{ h1}
    \KeywordTok{apply}\NormalTok{ h3}
    \KeywordTok{done}
\NormalTok{  dsimp }\KeywordTok{at}\NormalTok{ h4}
  \KeywordTok{linarith}
  \KeywordTok{done}
\end{Highlighting}
\end{Shaded}

\end{inpt}

\begin{outpt}

\begin{Shaded}
\begin{Highlighting}[]
\SpecialStringTok{No}\InformationTok{ }\SpecialStringTok{goals}
\end{Highlighting}
\end{Shaded}

\end{outpt}

\hypertarget{absolute-convergence}{%
\section{Absolute Convergence}\label{absolute-convergence}}

I will now display how I used to lean to demonstrate that a sequence
converges if and only if the absolute value of that sequence converges.

Laying out the goal

\begin{inpt}

\begin{Shaded}
\begin{Highlighting}[]
\KeywordTok{example}\NormalTok{ (s1 : ℕ → ℝ) : }
\NormalTok{    ConvergesTo s1 (0 : ℝ) ↔}
\NormalTok{    ConvergesTo (abs s1) (0 : ℝ) }
\NormalTok{    := }\KeywordTok{by}

  \SpecialCharTok{**}\WarningTok{done}\SpecialCharTok{::}
\end{Highlighting}
\end{Shaded}

\end{inpt}

\begin{outpt}

\begin{Shaded}
\begin{Highlighting}[]
\InformationTok{s1}\NormalTok{: ℕ → ℝ}
\NormalTok{⊢ ConvergesTo s1 0 }
\SpecialCharTok{\textgreater{}\textgreater{}}\NormalTok{  ↔ ConvergesTo |s1| 0}
\end{Highlighting}
\end{Shaded}

\end{outpt}

Step 1

\begin{inpt}

\begin{Shaded}
\begin{Highlighting}[]
\KeywordTok{example}\NormalTok{ (s1 : ℕ → ℝ) : }
\NormalTok{    ConvergesTo s1 (0 : ℝ) ↔}
\NormalTok{    ConvergesTo (abs s1) (0 : ℝ) }
\NormalTok{    := }\KeywordTok{by}
  \KeywordTok{rw}\NormalTok{ [ConvergesTo]}
  \KeywordTok{rw}\NormalTok{ [ConvergesTo]}

  \SpecialCharTok{**}\WarningTok{done}\SpecialCharTok{::}
\end{Highlighting}
\end{Shaded}

\end{inpt}

\begin{outpt}

\begin{Shaded}
\begin{Highlighting}[]
\InformationTok{s1}\NormalTok{: ℕ → ℝ}
\NormalTok{⊢ (∀ (ε : ℝ), ε \textgreater{} 0 → }
\SpecialCharTok{\textgreater{}\textgreater{}}\NormalTok{    ∃ N, ∀ (n : ℕ), }
\SpecialCharTok{\textgreater{}\textgreater{}}\NormalTok{    n ≥ N → }
\SpecialCharTok{\textgreater{}\textgreater{}}\NormalTok{    |s1 n {-} 0| \textless{} ε) ↔}
\SpecialCharTok{\textgreater{}\textgreater{}}\NormalTok{    ∀ (ε : ℝ), ε \textgreater{} 0 → }
\SpecialCharTok{\textgreater{}\textgreater{}}\NormalTok{    ∃ N, ∀ (n : ℕ), }
\SpecialCharTok{\textgreater{}\textgreater{}}\NormalTok{    n ≥ N → }
\SpecialCharTok{\textgreater{}\textgreater{}}\NormalTok{    |abs s1 n {-} 0| \textless{} ε}
\end{Highlighting}
\end{Shaded}

\end{outpt}

Step 2

\begin{inpt}

\begin{Shaded}
\begin{Highlighting}[]
\KeywordTok{example}\NormalTok{ (s1 : ℕ → ℝ) : }
\NormalTok{    ConvergesTo s1 (0 : ℝ) ↔}
\NormalTok{    ConvergesTo (abs s1) (0 : ℝ) }
\NormalTok{    := }\KeywordTok{by}
  \KeywordTok{rw}\NormalTok{ [ConvergesTo]}
  \KeywordTok{rw}\NormalTok{ [ConvergesTo]}
  \KeywordTok{have}\NormalTok{ h3 (x : ℕ) : |s1 x| = }
\NormalTok{      |abs s1 x| := }\KeywordTok{by}
    
    \KeywordTok{done}

  \SpecialCharTok{**}\WarningTok{done}\SpecialCharTok{::}
\end{Highlighting}
\end{Shaded}

\end{inpt}

\begin{outpt}

\begin{Shaded}
\begin{Highlighting}[]
\InformationTok{s1}\NormalTok{: ℕ → ℝ}
\InformationTok{x}\NormalTok{: ℕ}
\NormalTok{⊢ |s1 x| = |abs s1 x|}
\end{Highlighting}
\end{Shaded}

\end{outpt}

Step 3

\begin{inpt}

\begin{Shaded}
\begin{Highlighting}[]
\KeywordTok{example}\NormalTok{ (s1 : ℕ → ℝ) : }
\NormalTok{    ConvergesTo s1 (0 : ℝ) ↔}
\NormalTok{    ConvergesTo (abs s1) (0 : ℝ) }
\NormalTok{    := }\KeywordTok{by}
  \KeywordTok{rw}\NormalTok{ [ConvergesTo]}
  \KeywordTok{rw}\NormalTok{ [ConvergesTo]}
  \KeywordTok{have}\NormalTok{ h3 (x : ℕ) : |s1 x| = }
\NormalTok{      |abs s1 x| := }\KeywordTok{by}
    \KeywordTok{simp}\NormalTok{ [abs]}
    \KeywordTok{done}

  \SpecialCharTok{**}\WarningTok{done}\SpecialCharTok{::}
\end{Highlighting}
\end{Shaded}

\end{inpt}

\begin{outpt}

\begin{Shaded}
\begin{Highlighting}[]
\InformationTok{s1}\NormalTok{: ℕ → ℝ}
\InformationTok{h3}\NormalTok{: ∀ (x : ℕ), |s1 x| = }
\SpecialCharTok{\textgreater{}\textgreater{}}\NormalTok{  |abs s1 x|}
\NormalTok{⊢ (∀ (ε : ℝ), ε \textgreater{} 0 → }
\SpecialCharTok{\textgreater{}\textgreater{}}\NormalTok{    ∃ N, ∀ (n : ℕ), }
\SpecialCharTok{\textgreater{}\textgreater{}}\NormalTok{    n ≥ N → }
\SpecialCharTok{\textgreater{}\textgreater{}}\NormalTok{    |s1 n {-} 0| \textless{} ε) ↔}
\SpecialCharTok{\textgreater{}\textgreater{}}\NormalTok{    ∀ (ε : ℝ), ε \textgreater{} 0 → }
\SpecialCharTok{\textgreater{}\textgreater{}}\NormalTok{    ∃ N, ∀ (n : ℕ), }
\SpecialCharTok{\textgreater{}\textgreater{}}\NormalTok{    n ≥ N → }
\SpecialCharTok{\textgreater{}\textgreater{}}\NormalTok{    |abs s1 n {-} 0| \textless{} ε}
\end{Highlighting}
\end{Shaded}

\end{outpt}

Step 4

\begin{inpt}

\begin{Shaded}
\begin{Highlighting}[]
\KeywordTok{example}\NormalTok{ (s1 : ℕ → ℝ) : }
\NormalTok{    ConvergesTo s1 (0 : ℝ) ↔}
\NormalTok{    ConvergesTo (abs s1) (0 : ℝ) }
\NormalTok{    := }\KeywordTok{by}
  \KeywordTok{rw}\NormalTok{ [ConvergesTo]}
  \KeywordTok{rw}\NormalTok{ [ConvergesTo]}
  \KeywordTok{have}\NormalTok{ h3 (x : ℕ) : |s1 x| = }
\NormalTok{      |abs s1 x| := }\KeywordTok{by}
    \KeywordTok{simp}\NormalTok{ [abs]}
    \KeywordTok{done}
  \KeywordTok{apply}\NormalTok{ Iff.intro}
\NormalTok{  · }\CommentTok{{-}{-}Forwards}

\NormalTok{  · }\CommentTok{{-}{-}Reverse}

  \SpecialCharTok{**}\WarningTok{done}\SpecialCharTok{::}
\end{Highlighting}
\end{Shaded}

\end{inpt}

\begin{outpt}

\begin{Shaded}
\begin{Highlighting}[]
\InformationTok{s1}\NormalTok{: ℕ → ℝ}
\InformationTok{h3}\NormalTok{: ∀ (x : ℕ), |s1 x| = }
\SpecialCharTok{\textgreater{}\textgreater{}}\NormalTok{  |abs s1 x|}
\NormalTok{⊢ (∀ (ε : ℝ), ε \textgreater{} 0 → }
\SpecialCharTok{\textgreater{}\textgreater{}}\NormalTok{    ∃ N, ∀ (n : ℕ), }
\SpecialCharTok{\textgreater{}\textgreater{}}\NormalTok{    n ≥ N → }
\SpecialCharTok{\textgreater{}\textgreater{}}\NormalTok{    |s1 n {-} 0| \textless{} ε) →}
\SpecialCharTok{\textgreater{}\textgreater{}}\NormalTok{    ∀ (ε : ℝ), ε \textgreater{} 0 → }
\SpecialCharTok{\textgreater{}\textgreater{}}\NormalTok{    ∃ N, ∀ (n : ℕ), }
\SpecialCharTok{\textgreater{}\textgreater{}}\NormalTok{    n ≥ N → }
\SpecialCharTok{\textgreater{}\textgreater{}}\NormalTok{    |abs s1 n {-} 0| \textless{} ε}
\end{Highlighting}
\end{Shaded}

\end{outpt}

Step 5

\begin{inpt}

\begin{Shaded}
\begin{Highlighting}[]
\KeywordTok{example}\NormalTok{ (s1 : ℕ → ℝ) : }
\NormalTok{    ConvergesTo s1 (0 : ℝ) ↔}
\NormalTok{    ConvergesTo (abs s1) (0 : ℝ) }
\NormalTok{    := }\KeywordTok{by}
  \KeywordTok{rw}\NormalTok{ [ConvergesTo]}
  \KeywordTok{rw}\NormalTok{ [ConvergesTo]}
  \KeywordTok{have}\NormalTok{ h3 (x : ℕ) : |s1 x| = }
\NormalTok{      |abs s1 x| := }\KeywordTok{by}
    \KeywordTok{simp}\NormalTok{ [abs]}
    \KeywordTok{done}
  \KeywordTok{apply}\NormalTok{ Iff.intro}
\NormalTok{  · }\CommentTok{{-}{-}Forwards}
\NormalTok{    intro h1}

\NormalTok{  · }\CommentTok{{-}{-}Reverse}

  \SpecialCharTok{**}\WarningTok{done}\SpecialCharTok{::}
\end{Highlighting}
\end{Shaded}

\end{inpt}

\begin{outpt}

\begin{Shaded}
\begin{Highlighting}[]
\InformationTok{s1}\NormalTok{: ℕ → ℝ}
\InformationTok{h3}\NormalTok{: ∀ (x : ℕ), |s1 x| = }
\SpecialCharTok{\textgreater{}\textgreater{}}\NormalTok{  |abs s1 x|}
\InformationTok{h1}\NormalTok{: ∀ (ε : ℝ), ε \textgreater{} 0 → }
\SpecialCharTok{\textgreater{}\textgreater{}}\NormalTok{  ∃ N, ∀ (n : ℕ), }
\SpecialCharTok{\textgreater{}\textgreater{}}\NormalTok{  n ≥ N → |s1 n {-} 0| \textless{} ε}
\NormalTok{⊢ ∀ (ε : ℝ), ε \textgreater{} 0 → }
\SpecialCharTok{\textgreater{}\textgreater{}}\NormalTok{    ∃ N, ∀ (n : ℕ), }
\SpecialCharTok{\textgreater{}\textgreater{}}\NormalTok{    n ≥ N → }
\SpecialCharTok{\textgreater{}\textgreater{}}\NormalTok{    |abs s1 n {-} 0| \textless{} ε}
\end{Highlighting}
\end{Shaded}

\end{outpt}

Step 6

\begin{inpt}

\begin{Shaded}
\begin{Highlighting}[]
\KeywordTok{example}\NormalTok{ (s1 : ℕ → ℝ) : }
\NormalTok{    ConvergesTo s1 (0 : ℝ) ↔}
\NormalTok{    ConvergesTo (abs s1) (0 : ℝ) }
\NormalTok{    := }\KeywordTok{by}
  \KeywordTok{rw}\NormalTok{ [ConvergesTo]}
  \KeywordTok{rw}\NormalTok{ [ConvergesTo]}
  \KeywordTok{have}\NormalTok{ h3 (x : ℕ) : |s1 x| = }
\NormalTok{      |abs s1 x| := }\KeywordTok{by}
    \KeywordTok{simp}\NormalTok{ [abs]}
    \KeywordTok{done}
  \KeywordTok{apply}\NormalTok{ Iff.intro}
\NormalTok{  · }\CommentTok{{-}{-}Forwards}
\NormalTok{    intro h1}
    \KeywordTok{simp}

\NormalTok{  · }\CommentTok{{-}{-}Reverse}

  \SpecialCharTok{**}\WarningTok{done}\SpecialCharTok{::}
\end{Highlighting}
\end{Shaded}

\end{inpt}

\begin{outpt}

\begin{Shaded}
\begin{Highlighting}[]
\InformationTok{s1}\NormalTok{: ℕ → ℝ}
\InformationTok{h3}\NormalTok{: ∀ (x : ℕ), |s1 x| = }
\SpecialCharTok{\textgreater{}\textgreater{}}\NormalTok{  |abs s1 x|}
\InformationTok{h1}\NormalTok{: ∀ (ε : ℝ), ε \textgreater{} 0 → }
\SpecialCharTok{\textgreater{}\textgreater{}}\NormalTok{  ∃ N, ∀ (n : ℕ), }
\SpecialCharTok{\textgreater{}\textgreater{}}\NormalTok{  n ≥ N → |s1 n {-} 0| \textless{} ε}
\NormalTok{⊢ ∀ (ε : ℝ), 0 \textless{} ε → }
\SpecialCharTok{\textgreater{}\textgreater{}}\NormalTok{    ∃ N, ∀ (n : ℕ), }
\SpecialCharTok{\textgreater{}\textgreater{}}\NormalTok{    N ≤ n → }
\SpecialCharTok{\textgreater{}\textgreater{}}\NormalTok{    |abs s1 n| \textless{} ε}
\end{Highlighting}
\end{Shaded}

\end{outpt}

Step 7

\begin{inpt}

\begin{Shaded}
\begin{Highlighting}[]
\KeywordTok{example}\NormalTok{ (s1 : ℕ → ℝ) : }
\NormalTok{    ConvergesTo s1 (0 : ℝ) ↔}
\NormalTok{    ConvergesTo (abs s1) (0 : ℝ) }
\NormalTok{    := }\KeywordTok{by}
  \KeywordTok{rw}\NormalTok{ [ConvergesTo]}
  \KeywordTok{rw}\NormalTok{ [ConvergesTo]}
  \KeywordTok{have}\NormalTok{ h3 (x : ℕ) : |s1 x| = }
\NormalTok{      |abs s1 x| := }\KeywordTok{by}
    \KeywordTok{simp}\NormalTok{ [abs]}
    \KeywordTok{done}
  \KeywordTok{apply}\NormalTok{ Iff.intro}
\NormalTok{  · }\CommentTok{{-}{-}Forwards}
\NormalTok{    intro h1}
    \KeywordTok{simp}
    \KeywordTok{simp} \KeywordTok{at}\NormalTok{ h1}

\NormalTok{  · }\CommentTok{{-}{-}Reverse}

  \SpecialCharTok{**}\WarningTok{done}\SpecialCharTok{::}
\end{Highlighting}
\end{Shaded}

\end{inpt}

\begin{outpt}

\begin{Shaded}
\begin{Highlighting}[]
\InformationTok{s1}\NormalTok{: ℕ → ℝ}
\InformationTok{h3}\NormalTok{: ∀ (x : ℕ), |s1 x| = }
\SpecialCharTok{\textgreater{}\textgreater{}}\NormalTok{  |abs s1 x|}
\InformationTok{h1}\NormalTok{: ∀ (ε : ℝ), 0 \textless{} ε → }
\SpecialCharTok{\textgreater{}\textgreater{}}\NormalTok{  ∃ N, ∀ (n : ℕ), }
\SpecialCharTok{\textgreater{}\textgreater{}}\NormalTok{  N ≤ n → |s1 n| \textless{} ε}
\NormalTok{⊢ ∀ (ε : ℝ), 0 \textless{} ε → }
\SpecialCharTok{\textgreater{}\textgreater{}}\NormalTok{    ∃ N, ∀ (n : ℕ), }
\SpecialCharTok{\textgreater{}\textgreater{}}\NormalTok{    N ≤ n → }
\SpecialCharTok{\textgreater{}\textgreater{}}\NormalTok{    |abs s1 n| \textless{} ε}
\end{Highlighting}
\end{Shaded}

\end{outpt}

Step 8

\begin{inpt}

\begin{Shaded}
\begin{Highlighting}[]
\KeywordTok{example}\NormalTok{ (s1 : ℕ → ℝ) : }
\NormalTok{    ConvergesTo s1 (0 : ℝ) ↔}
\NormalTok{    ConvergesTo (abs s1) (0 : ℝ) }
\NormalTok{    := }\KeywordTok{by}
  \KeywordTok{rw}\NormalTok{ [ConvergesTo]}
  \KeywordTok{rw}\NormalTok{ [ConvergesTo]}
  \KeywordTok{have}\NormalTok{ h3 (x : ℕ) : |s1 x| = }
\NormalTok{      |abs s1 x| := }\KeywordTok{by}
    \KeywordTok{simp}\NormalTok{ [abs]}
    \KeywordTok{done}
  \KeywordTok{apply}\NormalTok{ Iff.intro}
\NormalTok{  · }\CommentTok{{-}{-}Forwards}
\NormalTok{    intro h1}
    \KeywordTok{simp}
    \KeywordTok{simp} \KeywordTok{at}\NormalTok{ h1}
    \KeywordTok{simp}\NormalTok{ [← h3]}

\NormalTok{  · }\CommentTok{{-}{-}Reverse}

  \SpecialCharTok{**}\WarningTok{done}\SpecialCharTok{::}
\end{Highlighting}
\end{Shaded}

\end{inpt}

\begin{outpt}

\begin{Shaded}
\begin{Highlighting}[]
\InformationTok{s1}\NormalTok{: ℕ → ℝ}
\InformationTok{h3}\NormalTok{: ∀ (x : ℕ), |s1 x| = }
\SpecialCharTok{\textgreater{}\textgreater{}}\NormalTok{  |abs s1 x|}
\InformationTok{h1}\NormalTok{: ∀ (ε : ℝ), 0 \textless{} ε → }
\SpecialCharTok{\textgreater{}\textgreater{}}\NormalTok{  ∃ N, ∀ (n : ℕ), }
\SpecialCharTok{\textgreater{}\textgreater{}}\NormalTok{  N ≤ n → |s1 n| \textless{} ε}
\NormalTok{⊢ ∀ (ε : ℝ), 0 \textless{} ε → }
\SpecialCharTok{\textgreater{}\textgreater{}}\NormalTok{    ∃ N, ∀ (n : ℕ), }
\SpecialCharTok{\textgreater{}\textgreater{}}\NormalTok{    N ≤ n → }
\SpecialCharTok{\textgreater{}\textgreater{}}\NormalTok{    |s1 n| \textless{} ε}
\end{Highlighting}
\end{Shaded}

\end{outpt}

Step 9

\begin{inpt}

\begin{Shaded}
\begin{Highlighting}[]
\KeywordTok{example}\NormalTok{ (s1 : ℕ → ℝ) : }
\NormalTok{    ConvergesTo s1 (0 : ℝ) ↔}
\NormalTok{    ConvergesTo (abs s1) (0 : ℝ) }
\NormalTok{    := }\KeywordTok{by}
  \KeywordTok{rw}\NormalTok{ [ConvergesTo]}
  \KeywordTok{rw}\NormalTok{ [ConvergesTo]}
  \KeywordTok{have}\NormalTok{ h3 (x : ℕ) : |s1 x| = }
\NormalTok{      |abs s1 x| := }\KeywordTok{by}
    \KeywordTok{simp}\NormalTok{ [abs]}
    \KeywordTok{done}
  \KeywordTok{apply}\NormalTok{ Iff.intro}
\NormalTok{  · }\CommentTok{{-}{-}Forwards}
\NormalTok{    intro h1}
    \KeywordTok{simp}
    \KeywordTok{simp} \KeywordTok{at}\NormalTok{ h1}
    \KeywordTok{simp}\NormalTok{ [← h3]}
    \KeywordTok{apply}\NormalTok{ h1}
\NormalTok{  · }\CommentTok{{-}{-}Reverse}

  \SpecialCharTok{**}\WarningTok{done}\SpecialCharTok{::}
\end{Highlighting}
\end{Shaded}

\end{inpt}

\begin{outpt}

\begin{Shaded}
\begin{Highlighting}[]
\InformationTok{s1}\NormalTok{: ℕ → ℝ}
\InformationTok{h3}\NormalTok{: ∀ (x : ℕ), |s1 x| = }
\SpecialCharTok{\textgreater{}\textgreater{}}\NormalTok{  |abs s1 x|}
\NormalTok{⊢ (∀ (ε : ℝ), ε \textgreater{} 0 → }
\SpecialCharTok{\textgreater{}\textgreater{}}\NormalTok{    ∃ N, ∀ (n : ℕ), }
\SpecialCharTok{\textgreater{}\textgreater{}}\NormalTok{    n ≥ N → }
\SpecialCharTok{\textgreater{}\textgreater{}}\NormalTok{    |abs s1 n {-} 0| \textless{} ε) →}
\SpecialCharTok{\textgreater{}\textgreater{}}\NormalTok{    ∀ (ε : ℝ), ε \textgreater{} 0 → }
\SpecialCharTok{\textgreater{}\textgreater{}}\NormalTok{    ∃ N, ∀ (n : ℕ), }
\SpecialCharTok{\textgreater{}\textgreater{}}\NormalTok{    n ≥ N → }
\SpecialCharTok{\textgreater{}\textgreater{}}\NormalTok{    |s1 n {-} 0| \textless{} ε}
\end{Highlighting}
\end{Shaded}

\end{outpt}

Step 10

\begin{inpt}

\begin{Shaded}
\begin{Highlighting}[]
\KeywordTok{example}\NormalTok{ (s1 : ℕ → ℝ) : }
\NormalTok{    ConvergesTo s1 (0 : ℝ) ↔}
\NormalTok{    ConvergesTo (abs s1) (0 : ℝ) }
\NormalTok{    := }\KeywordTok{by}
  \KeywordTok{rw}\NormalTok{ [ConvergesTo]}
  \KeywordTok{rw}\NormalTok{ [ConvergesTo]}
  \KeywordTok{have}\NormalTok{ h3 (x : ℕ) : |s1 x| = }
\NormalTok{      |abs s1 x| := }\KeywordTok{by}
    \KeywordTok{simp}\NormalTok{ [abs]}
    \KeywordTok{done}
  \KeywordTok{apply}\NormalTok{ Iff.intro}
\NormalTok{  · }\CommentTok{{-}{-}Forwards}
\NormalTok{    intro h1}
    \KeywordTok{simp}
    \KeywordTok{simp} \KeywordTok{at}\NormalTok{ h1}
    \KeywordTok{simp}\NormalTok{ [← h3]}
    \KeywordTok{apply}\NormalTok{ h1}
\NormalTok{  · }\CommentTok{{-}{-}Reverse}
\NormalTok{    intro h1}
    \KeywordTok{simp}
    \KeywordTok{simp} \KeywordTok{at}\NormalTok{ h1}
\NormalTok{    simp\_rw [h3]}

  \SpecialCharTok{**}\WarningTok{done}\SpecialCharTok{::}
\end{Highlighting}
\end{Shaded}

\end{inpt}

\begin{outpt}

\begin{Shaded}
\begin{Highlighting}[]
\InformationTok{s1}\NormalTok{: ℕ → ℝ}
\InformationTok{h3}\NormalTok{: ∀ (x : ℕ), |s1 x| = }
\SpecialCharTok{\textgreater{}\textgreater{}}\NormalTok{  |abs s1 x|}
\InformationTok{h1}\NormalTok{: ∀ (ε : ℝ), 0 \textless{} ε → }
\SpecialCharTok{\textgreater{}\textgreater{}}\NormalTok{  ∃ N, ∀ (n : ℕ), }
\SpecialCharTok{\textgreater{}\textgreater{}}\NormalTok{  N ≤ n → |abs s1 n| \textless{} ε}
\NormalTok{⊢ ∀ (ε : ℝ), 0 \textless{} ε → }
\SpecialCharTok{\textgreater{}\textgreater{}}\NormalTok{    ∃ N, ∀ (n : ℕ),}
\SpecialCharTok{\textgreater{}\textgreater{}}\NormalTok{    N ≤ n → }
\SpecialCharTok{\textgreater{}\textgreater{}}\NormalTok{    |abs s1 n| \textless{} ε}
\end{Highlighting}
\end{Shaded}

\end{outpt}

Step 11

\begin{inpt}

\begin{Shaded}
\begin{Highlighting}[]
\KeywordTok{example}\NormalTok{ (s1 : ℕ → ℝ) : }
\NormalTok{    ConvergesTo s1 (0 : ℝ) ↔}
\NormalTok{    ConvergesTo (abs s1) (0 : ℝ) }
\NormalTok{    := }\KeywordTok{by}
  \KeywordTok{rw}\NormalTok{ [ConvergesTo]}
  \KeywordTok{rw}\NormalTok{ [ConvergesTo]}
  \KeywordTok{have}\NormalTok{ h3 (x : ℕ) : |s1 x| = }
\NormalTok{      |abs s1 x| := }\KeywordTok{by}
    \KeywordTok{simp}\NormalTok{ [abs]}
    \KeywordTok{done}
  \KeywordTok{apply}\NormalTok{ Iff.intro}
\NormalTok{  · }\CommentTok{{-}{-}Forwards}
\NormalTok{    intro h1}
    \KeywordTok{simp}
    \KeywordTok{simp} \KeywordTok{at}\NormalTok{ h1}
    \KeywordTok{simp}\NormalTok{ [← h3]}
    \KeywordTok{apply}\NormalTok{ h1}
\NormalTok{  · }\CommentTok{{-}{-}Reverse}
\NormalTok{    intro h1}
    \KeywordTok{simp}
    \KeywordTok{simp} \KeywordTok{at}\NormalTok{ h1}
\NormalTok{    simp\_rw [h3]}
    \KeywordTok{apply}\NormalTok{ h1}
  \KeywordTok{done}
\end{Highlighting}
\end{Shaded}

\end{inpt}

\begin{outpt}

\begin{Shaded}
\begin{Highlighting}[]
\SpecialStringTok{No}\InformationTok{ }\SpecialStringTok{goals}
\end{Highlighting}
\end{Shaded}

\end{outpt}

\hypertarget{convergence-of-a-specific-sequence}{%
\section{Convergence of a Specific
Sequence}\label{convergence-of-a-specific-sequence}}

example of lean bad.

\begin{Shaded}
\begin{Highlighting}[]
\KeywordTok{example}\NormalTok{ : ConvergesTo (fun (n : ℕ) ↦ }
\NormalTok{    ((2 * n) / (n + 1))) 2 := }\KeywordTok{by}
\NormalTok{  intro ε}
\NormalTok{  intro h1}
  \KeywordTok{obtain}\NormalTok{ ⟨k, h13⟩ := }
\NormalTok{    exists\_nat\_gt (2 / ε {-} 1) }\CommentTok{{-}{-}Archimedean Property}
\NormalTok{  use k}
\NormalTok{  intro n}
\NormalTok{  intro h2}
\NormalTok{  dsimp}
  \KeywordTok{have}\NormalTok{ h3 : (2 : ℝ) = 2 * ((n + 1) / (n + 1)) := }\KeywordTok{by}
    \KeywordTok{have}\NormalTok{ h4 : ((n + 1) / (n + 1)) = }
\NormalTok{        (n + 1) * ((n + 1) : ℝ)⁻¹ := }\KeywordTok{by}
      \KeywordTok{rfl}
      \KeywordTok{done}
    \KeywordTok{rw}\NormalTok{ [h4]}
    \KeywordTok{have}\NormalTok{ h5 : (n + 1) * ((n + 1) : ℝ)⁻¹ = 1 := }\KeywordTok{by}
      \KeywordTok{rw}\NormalTok{ [mul\_inv\_cancel]}
      \KeywordTok{exact}\NormalTok{ Nat.cast\_add\_one\_ne\_zero n}
      \KeywordTok{done}
    \KeywordTok{rw}\NormalTok{ [h5]}
    \KeywordTok{exact}\NormalTok{ Eq.symm (mul\_one 2)}
    \KeywordTok{done}
\NormalTok{  nth\_rewrite 2 [h3]}
  \KeywordTok{have}\NormalTok{ h6 : 2 * ((↑n + 1) : ℝ) / (↑n + 1) = }
\NormalTok{      ((2 * n) + 2) / (n + 1) := }\KeywordTok{by}
    \KeywordTok{rw}\NormalTok{ [Distribute n]}
    \KeywordTok{done}
  \KeywordTok{have}\NormalTok{ h7 : 2 * (((↑n + 1) : ℝ) / (↑n + 1)) = }
\NormalTok{      2 * (↑n + 1) / (↑n + 1) := }\KeywordTok{by}
    \KeywordTok{rw}\NormalTok{ [← mul\_div\_assoc 2 ((n + 1) : ℝ) ((n + 1) : ℝ)]}
    \KeywordTok{done}
  \KeywordTok{rw}\NormalTok{ [h7]}
  \KeywordTok{rw}\NormalTok{ [h6]}
  \KeywordTok{rw}\NormalTok{ [div\_sub\_div\_same (2 * n : ℝ) (2 * n + 2) (n + 1)]}
  \KeywordTok{rw}\NormalTok{ [sub\_add\_cancel\textquotesingle{}]}
  \KeywordTok{rw}\NormalTok{ [abs\_div]}
  \KeywordTok{simp}
  \KeywordTok{have}\NormalTok{ h8 : |(↑n + 1 : ℝ)| = ↑n + 1 := }\KeywordTok{by}
    \KeywordTok{simp}
    \KeywordTok{apply}\NormalTok{ LT.lt.le (Nat.cast\_add\_one\_pos ↑n)}
    \KeywordTok{done}
  \KeywordTok{rw}\NormalTok{ [h8]}
  \KeywordTok{have}\NormalTok{ h9 : (2 : ℝ) / (↑n + 1) ≤ 2 / (k + 1) := }\KeywordTok{by}
    \KeywordTok{apply}\NormalTok{ div\_le\_div\_of\_le\_left}
\NormalTok{    · }\CommentTok{{-}{-}case 1}
      \KeywordTok{linarith}
      \KeywordTok{done}
\NormalTok{    · }\CommentTok{{-}{-}case 2}
      \KeywordTok{exact}\NormalTok{ Nat.cast\_add\_one\_pos k}
      \KeywordTok{done}
\NormalTok{    · }\CommentTok{{-}{-}case 3}
\NormalTok{      convert add\_le\_add\_right h2 1}
      \KeywordTok{apply}\NormalTok{ Iff.intro}
\NormalTok{      · }\CommentTok{{-}{-}subcase 1}
        \KeywordTok{exact}\NormalTok{ fun a =\textgreater{} Nat.add\_le\_add\_right h2 1}
        \KeywordTok{done}
\NormalTok{      · }\CommentTok{{-}{-}subcase 2}
\NormalTok{        intro h14}
        \KeywordTok{apply}\NormalTok{ add\_le\_add\_right}
        \KeywordTok{exact}\NormalTok{ Iff.mpr Nat.cast\_le h2}
        \KeywordTok{done}
      \KeywordTok{done}
    \KeywordTok{done}
  \KeywordTok{have}\NormalTok{ h10 : 2 / (k + 1) \textless{} 2 / (2 / ε {-} 1 + 1) := }\KeywordTok{by}
    \KeywordTok{apply}\NormalTok{ div\_lt\_div\_of\_lt\_left}
\NormalTok{    · }\CommentTok{{-}{-}case 1}
      \KeywordTok{linarith}
      \KeywordTok{done}
\NormalTok{    · }\CommentTok{{-}{-}case 2}
      \KeywordTok{simp}
      \KeywordTok{apply}\NormalTok{ div\_pos}
      \KeywordTok{linarith}
      \KeywordTok{apply}\NormalTok{ h1}
      \KeywordTok{done}
\NormalTok{    · }\CommentTok{{-}{-}case 3}
\NormalTok{      convert add\_le\_add\_right h2 1}
      \KeywordTok{apply}\NormalTok{ Iff.intro}
\NormalTok{      · }\CommentTok{{-}{-}subcase 1}
\NormalTok{        intro h11}
        \KeywordTok{exact}\NormalTok{ Nat.add\_le\_add\_right h2 1}
        \KeywordTok{done}
\NormalTok{      · }\CommentTok{{-}{-}subcase 2}
\NormalTok{        intro h11}
        \KeywordTok{have}\NormalTok{ h12 : 2 / ε {-} 1 \textless{} (k : ℝ) := }\KeywordTok{by}
          \KeywordTok{simp}\NormalTok{ only []}
          \KeywordTok{apply}\NormalTok{ h13}
          \KeywordTok{done}
        \KeywordTok{exact}\NormalTok{ add\_lt\_add\_right h12 1}
        \KeywordTok{done}
      \KeywordTok{done}
    \KeywordTok{done}
  \KeywordTok{calc}
\NormalTok{    2 / (↑n + 1) ≤ (2 : ℝ) / (k + 1) := }\KeywordTok{by}
      \KeywordTok{apply}\NormalTok{ h9}
      \KeywordTok{done}
\NormalTok{    \_ \textless{} (2 : ℝ) / (2 / ε {-} 1 + 1) := }\KeywordTok{by}
      \KeywordTok{apply}\NormalTok{ h10}
      \KeywordTok{done}
\NormalTok{    \_ = ε := }\KeywordTok{by}
\NormalTok{      ring\_nf}
      \KeywordTok{apply}\NormalTok{ inv\_inv}
      \KeywordTok{done}
  \KeywordTok{done}
\end{Highlighting}
\end{Shaded}

\hypertarget{proofs-involving-quantifiers}{%
\section{3.3. Proofs Involving
Quantifiers}\label{proofs-involving-quantifiers}}

\begin{thm}
Suppose \(B\) is a set and \(\mathcal{F}\) is a family of sets. If
\(\bigcup\mathcal{F} \subseteq B\) then
\(\mathcal{F} \subseteq \mathscr{P}(B)\).

\end{thm}

\begin{proof}

Suppose \(\bigcup \mathcal{F} \subseteq B\). Let \(x\) be an arbitrary
element of \(\mathcal{F}\). Let \(y\) be an arbitrary element of \(x\).
Since \(y \in x\) and \(x \in \mathcal{F}\), by the definition of
\(\bigcup \mathcal{F}\), \(y \in \bigcup \mathcal{F}\). But then since
\(\bigcup \mathcal{F} \subseteq B\), \(y \in B\). Since \(y\) was an
arbitrary element of \(x\), we can conclude that \(x \subseteq B\), so
\(x \in \mathscr{P}(B)\). But \(x\) was an arbitrary element of
\(\mathcal{F}\), so this shows that
\(\mathcal{F} \subseteq \mathscr{P}(B)\), as required. \excl{~□}\qedhere

\end{proof}

\begin{nthm}{Theorem 3.4.7}
For every integer \(n\), \(6 \mid n\) iff \(2 \mid n\) and \(3 \mid n\).

\end{nthm}

\begin{proof}

Let \(n\) be an arbitrary integer.

(\(\to\)) Suppose \(6 \mid n\). Then we can choose an integer \(k\) such
that \(6k=n\). Therefore \(n = 6k = 2(3k)\), so \(2 \mid n\), and
similarly \(n = 6k = 3(2k)\), so \(3 \mid n\).

(\(\leftarrow\)) Suppose \(2 \mid n\) and \(3 \mid n\). Then we can
choose integers \(j\) and \(k\) such that \(n = 2j\) and \(n = 3k\).
Therefore \(6(j-k) = 6j - 6k = 3(2j) - 2(3k) = 3n - 2n = n\), so
\(6 \mid n\). \excl{~□}\qedhere

\end{proof}

\begin{mdsk}

\end{mdsk}

\bookmarksetup{startatroot}

\hypertarget{conclusions}{%
\chapter{Conclusions}\label{conclusions}}

\[
[x]_R = \{y \in A \mid yRx\}.
\] The set whose elements are all of these equivalence classes is called
\(A\) \emph{mod} \(R\). It is written \(A/R\), so \[
A/R = \{[x]_R \mid x \in A\}.
\] Note that \(A/R\) is a set whose elements are sets: for each
\(x \in A\), \([x]_R\) is a subset of \(A\), and \([x]_R \in A/R\).

\bookmarksetup{startatroot}

\hypertarget{works-cited}{%
\chapter{Works Cited}\label{works-cited}}

This work had been formatted and styled from the book \emph{How To Prove
It With Lean}, written by Daniel J. Velleman. \emph{How To Prove It With
Lean} contains short excerpts from \emph{How To Prove It: A Structured
Approach, 3rd Edition}, by Daniel J. Velleman and published by Cambridge
University Press.

\begin{Shaded}
\begin{Highlighting}[]
\KeywordTok{example}\NormalTok{ : square1 = square2 := }\KeywordTok{by} \KeywordTok{rfl}

\SpecialCharTok{++}\DocumentationTok{\#eval}\SpecialCharTok{::}\NormalTok{ square1 7     }\CommentTok{{-}{-}Answer: 49}
\end{Highlighting}
\end{Shaded}

\bookmarksetup{startatroot}

\hypertarget{additional-space}{%
\chapter{Additional space}\label{additional-space}}

Extra chapter to write more things if needed!!



\end{document}
